\section{Συζήτηση}

\subsection{Συζήτηση αποτελεσμάτων των μηχανικών μοντέλων}
Χρησιμοποιούμε τις μετρικές αξιολόγησης των μηχανικών μοντέλων για να δούμε ποιες με μια γρήγορα ματιά ποιες υποπεριπτώσεις έχουν βγάλει αξιόλογα αποτελέσματα.
Αλλά επειδή αυτές οι μετρικές δεν μας δίνουν πλήρη εικόνα πως πήγε η πρόβλεψη, χρησιμοποιούμε τα διαγράμματα για να δούμε πόσο κοντινές είναι οι προβλέψεις.

Κάτι που παρατηρείται άμεσα κοιτώντας τα διαγράμματα, είναι ότι πολλές περιπτώσεις τα τελικά αποτελέσματα ξεφεύγουν από τα όρια των τοίχων του δωματίου.
Λαμβάνοντας αυτό υπόψη, μπορούμε να δηλώσουμε ότι σε διάφορες υποπεριπτώσεις που ισχύει αυτό, οι μετρικές αξιολογήσεις έχουν χειρότερα αποτελέσματα από την πραγματικότητα,
αφού σε περιπτώσεις που μελετάμε τα αποτελέσματα για ένα δωμάτιο, βγάζει νόημα να θεωρήσουμε ότι η πηγή ρύπου είναι δίπλα στο τοίχο.
Θα μπορούσαμε να "φτιάξουμε" τα αποτελέσματα αυτά, φέρνοντας τα σημεία στα "νόμιμα" όρια του δωματίου, αλλά θα μας έδινε ψευδά αποτελέσματα σε περίπτωση που θέλαμε να μελετήσουμε πολλαπλά δωμάτια,
ή σε περιπτώσεις μεγάλων δωματίων, τα αποτελέσματα θα υποδεικνυόντουσαν χειρότερα, αφού η απόκλιση των βασικών μοντέλων να βρουν την ευκλείδεια απόσταση θα γινόταν φανερή.
 
Μπορούμε να δούμε ότι όταν παίρνουμε όλα τα δεδομένα, όπου τα αποτελέσματα είναι στο \autoref{tab:all_data},τα προβλεπόμενα σημεία συγκεντρώνονται σε ένα μέρος του δωματίου και όχι σε όλο το εύρος τους, όπως θα έπρεπε.
Παρατηρήθηκε ότι όσο πιο "κακό" ήταν ένα βασικό μηχανικό μοντέλο, οι προβλέψεις για την ευκλείδεια απόσταση κυμαίνονταν προς μία τιμή, συνήθως την πιο μεσαία τιμή, και όχι σε όλο το εύρος ευκλείδιων αποστάσεων.
Μπορούμε να δούμε ότι και τα τελικά αποτελέσματα, όπως και τα βασικά μοντέλα, δεν υπάρχει κάπου R\textsuperscript{2} μεγαλύτερο του 0.5.
Οπότε με βάση τις προσεγγίσεις που ακολουθήσαμε, δεν μπορούμε να βρούμε την θέση πηγής ρύπου. Οι μέθοδοι διαχείρισης των outliers δεν αλλάζουν την εικόνα.

Πηγαίνοντας να εκμεταλλευτούμε την ιδιότητα της υπερευαισθησίας των MOS αισθητήρων αέρα που χρησιμοποιεί η BME680,όπου οδηγούν τα VOC να φτάνουν στην μέγιστη τιμή,
αναζητήσαμε προσεγγίσεις όπου εκμεταλλευόμαστε την μεγάλη διαφοροποίηση που έχουν οι πιο κοντινές πηγές ρύπου,
ειδικά οι πηγές ρύπου με ευκλείδεια απόσταση κάτω από μισό μέτρο, σε σχέση με τις υπόλοιπες πηγές ρύπου.

Γενικά, να σχολιάσουμε ότι για τις αποστάσεις που μιλάμε, δεν χρειάζεται να έχουμε πλήρη ταύτιση μεταξύ πραγματικής και προβλεπόμενης θέσης πηγής ρύπου.
Αν μπορούμε να εντοπίσει σταθερά την γενική υποπεριοχή που βρίσκεται μία πηγή ρύπου, μπορεί ο χρήστης από τα αντικείμενα που υπάρχουν στον άμεσο υποχώρο να εντοπίσει το αντικείμενο ή το σημείο από που εξάγεται ο ρύπος στον εσωτερικό αέρα.

Μπορούμε να δούμε ότι για την υποκατηγορία δεδομένων \textbf{Πιο κοντινή πηγή},που βρίσκεται στο \autoref{tab:closest_data},
ότι όλες οι προσεγγίσεις παρήγαγαν πολύ καλύτερα αποτελέσματα από θέμα επίδοσης.
Βεβαίως αν κοιτάξουμε τα διαγράμματα επιμέρους αποτελεσμάτων, θα δούμε ότι για τις περιπτώσεις \textbf{Όχι Διαχείριση Outliers} και \textbf{Διαχείριση Outliers σε συνολικό επίπεδο},
οι προβλεπόμενες πηγές δεν θέλουμε να είναι 

Τα καλύτερα αποτελέσματα προέρχονται από την υποπερίπτωση \textbf{Διαχείριση Outliers Ανά στήλη},όπου έχουμε R\textsuperscript{2} = 0.835465.
Μπορούμε να δούμε ότι ενώ δεν έχουμε εκατό τις εκατό ακριβή προσέγγιση των πηγών, έχουμε μέση ευκλείδεια απόσταση στο μισό μέτρο, όπου στην πραγματικότητα είναι μικρότερη,
αφού η πηγή ρύπου που φαίνεται στα διαγράμματα ως \textbf{Id=0} είναι εκτός του δωματίου/

Εδώ μπορούμε να δούμε την ευαισθησία που έχει ο απλός τριαπελευρικός εντοπισμός,
γιατί ο λόγος που δεν έχουμε πολύ καλή προσέγγιση, ενώ το R\textsuperscript{2} των 
βασικών μοντέλων των αισθητήρων με \textbf{Id=0} και \textbf{Id=1} είναι στο 0.9,το μοντέλο του αισθητήρα  με \textbf{Id=2} είχε R\textsuperscript{2} = 0.7393.

Μπορούμε να δούμε αντίστοιχα αποτελέσματα για την υποκατηγορία δεδομένων \textbf{Δύο πιο κοντινές πηγές ανά αισθητήρα},που βρίσκεται στο \autoref{tab:second_closest_data}.
Οι μετρικές είναι αρκετά καλές για \textbf{Διαχείριση Outliers σε συνολικό επίπεδο} και \textbf{Διαχείριση Outliers Ανά στήλη},όχι για \textbf{Όχι Διαχείριση Outliers},που βγάζει αρνητικό R\textsuperscript{2}.
Και εδώ, η \textbf{Διαχείριση Outliers Ανά στήλη},αποτελεί την καλύτερη επιλογή

Δεν είναι περίεργο που δεν καταφέραμε να δημιουργήσουμε μηχανικά μοντέλα που να βρίσκουν την θέση πηγή ρύπου σε όλο το δωμάτιο.
Οι χρονοσειρές του κάθε πειράματος στην ίδια θέση δεν είχαν κάποια κοινή συμπεριφορά, για να μπορέσει να παραχθεί κάποιο μηχανικό μοντέλο.
Τα δεδομένα που είχαμε δεν είναι πολλά, έχουμε 6 με 7 πειράματα ανά θέση, με σύνολο πειραμάτων που λάβαμε υπόψη να ανέρχονται στα 98 πειράματα. 

Αλλά δεν θα  είχε νόημα να κάτσουμε να κάνουμε περισσότερα πειράματα, αφού η μεγάλη μεταβλητότητα των δεδομένων θα σήμαινε ότι ήταν σίγουρο ότι δεν θα βοηθούσανε
ένα δύο πειράματα ανά θέση ακόμα.
Και λαμβάνοντας υπόψη ότι κάθε πείραμα έπαιρνε μισή ώρα περίπου και δοκιμάζαμε 16 θέσεις, συν την διαδικασία calibration και κάποια πειράματα που λαμβάναμε.
χωρίς πηγή ή δίπλα στους αισθητήρες, δεν άξιζε ο χρόνος που θα χρειαζόταν να περάσουμε για να μαζέψουμε τα δεδομένα, λαμβάνοντας υπόψη ότι μελετάμε ένα μόνο δωμάτιο σε σχετικά σταθερές συνθήκες:

Ο λόγος που τα δεδομένα δεν μπορούν να δημιουργήσουν μηχανικά μοντέλα είναι για διαφόρους λόγους:

\begin{enumerate}
    \item Τα εσωτερικά μοντέλα της BSEC λαμβάνουν υπόψη προηγούμενη δεδομένα για να βρουν τις τιμές των virtual εξόδους των τιμών που έχουν σχέση με την ποιότητα αέρα, στην περίπτωση μας τα VOC.
    Ακόμα και αν δεν κρατούσαμε δεδομένα στην μνήμη Flash, δεν θα μπορούσαμε να αναπαράγουμε τέλεια τα δεδομένα, ώστε να έχουν τα πειράματα με κοινή θέση πηγής ρύπου κοινό ιστορικό.
    Οπότε μιλάμε για χρονικό μεταβαλλόμενο σύστημα.
    \item Οι αισθητήρες έχουν δοκιμαστεί για να αναγνωρίζουν κυρίως breath VOCs,δηλαδή πτητικές χημικές ενώσεις που βρίσκονται στην ανθρώπινη ανάσα, οπότε και μόνο η μικρή παρουσία ενός ανθρώπου για να εισάγει την πηγή ρύπου επηρέαζε το σύστημα.
    \item Ακόμα και να μην ισχύει το παραπάνω, οι αισθητήρες εντοπίζουν TVOCs, δεν μπορούν να ξεχωρίσουν την αιθανόλη ξεχωριστά. Οπότε και τα TVOCs από τοίχους και έπιπλα μπορεί να επηρεάζουν το αποτέλεσμα.
    \item Διαφορετικές περιβαλλοντικές συνθήκες, ακόμα και αν είναι παρόμοιες, αφού έγιναν σε κοντινό χρονικό διάστημα όλες οι μετρήσεις.
    \item Εκ φύσεως περίπλοκη συμπεριφορά των αερίων.
    \item Αδυναμίες που έχουν γενικά οι LCS.
\end{enumerate}

Με βάση τους παραπάνω περιορισμούς, το ότι βρήκαμε περιπτώσεις με ικανοποιητικά αποτελέσματα είναι πολύ θετικές, παρόλο που μιλάμε για συγκεκριμένες υποπεριπτώσεις των πειραμάτων.
Οι BME680 αποδεικνύουν ότι μπορούν να χρησιμοποιηθούν για την καταμέτρηση της ποιότητας του αέρα, με την αντίδραση τους στα γεγονότα να είναι αρκούντως ικανοποιητική.
Δεν μπορούμε να ξέρουμε την ακρίβεια των μετρήσεων, αφού δεν χρησιμοποιήσαμε κάποιο όργανο αναφοράς.
Αλλά με τήρηση συνθηκών που προτείνονται άμεσα ή έμμεσα από την εταιρία, δηλαδή calibration που κρατάει τον αισθητήρα σε ρυπαρό αέρα για μισή ώρα ή συνεχή λειτουργία του αισθητήρα για μέρες, ώστε να έχει ιστορικό του δωματίου,
μπορεί να προσφέρει καλύτερα αποτελέσματα. 

Κάτι που πρέπει να σχολιάσουμε είναι ότι τα μηχανικά μοντέλα και η μέθοδος τριπλευρικού εντοπισμού δούλεψε σε συνθήκες που ξεφεύγουν τυπικών τακτικών διαχείρισης μηχανικών μοντέλων,
που είναι η κανονοικοποίηση (\texttt{normalize}) της κάθε γραμμής στην μονάδας της ευκλείδειας νόρμας, που ξεφεύγει από τον \texttt{StandardScaler} ή \texttt{MinMaxScaler} που χρησιμοποιείται.
Επίσης, κόβουμε τα Outliers μετά από την κανονικοποίηση, ενώ κανονικά γίνεται πριν.

Ένας λόγος που μπορούμε να αποδώσουμε την αποτελεσματικότητα της κανονικοποίησης σε σχέση με άλλα μετασχηματιστές κλίμακας,
είναι ότι δεν επηρεάζεται από τις ακραίες τιμές που υπάρχουν στις στήλες του πίνακα εισόδου, οι μεγάλες αυξήσεις των τιμών μπορούν να ξεχωρίσουν χωρίς να υπερισχύουν των υπολοίπων δεδομένων
και ο ακριβής αριθμός των τιμών δεν μας νοιάζουν τόσο.
Για την διαχείριση Outliers ανά στήλη, μπορούμε να συμπεράνουμε ότι με το να κόψουμε τις ακραίες τιμές όταν έχει αλλάξει η κλίμακα, εξασφαλίζουμε την διατήρηση σημαντικού βαθμού πληροφορίας.

Τέλος, πρέπει να αναφέρουμε ότι τα αποτελέσματα των μηχανικών μοντέλων βελτιώθηκαν πολύ, όταν αλλάξαμε τα δεδομένα που λαμβάνουμε υπόψη κατά το τέλος της.
Ξεκινήσαμε την ανάλυση των δεδομένων αναλύοντας από 30 δευτερόλεπτα πριν την εισαγωγή 
πηγής ρύπου, μέχρι 300 δευτερόλεπτα μετά από την εισαγωγή πηγής ρύπου, δηλαδή 330 δευτερόλεπτα.
Η αλλαγή στα 265 δευτερόλεπτα προέκυψε ως ανάγκη μιας προσέγγισης που ακολουθήσαμε, να βρεθεί αρκετός χώρος ολίσθησης μιας υποκατηγορίας των δεδομένων που είχε επιλεχθεί από τα δεδομένα που ανήκουν στα πειράματα που ανήκουν \textbf{Δύο πιο κοντινές πηγές ανά αισθητήρα},
χωρίς να ξεχωρίζουμε αν έχουμε κλειστή ή ανοιχτή την πόρτα.

Η προσπάθεια αυτή δεν έβγαλε αποτέλεσμα, αλλά χάρη στην αλλαγή αυτή, βελτιώθηκαν τα μοντέλα για τις υποπεριπτώσεις \textbf{Πιο κοντινή πηγή ανά αισθητήρα} και \textbf{Δύο πιο κοντινές πηγές ανά αισθητήρα}.
Με βάση αυτή την εμπειρία, μπορούμε να πούμε ότι φαίνεται ότι υπάρχουν βελτιώσεις που μπορούν να γίνουν και στο κομμάτι της επεξεργασίας των δεδομένων και στα μηχανικά μοντέλα, και δεν είναι μόνο ζήτημα αρχικών δεδομένων 

\subsection{Μελλοντικές κατευθύνσεις}

Η εργασία που τελέστηκε είχε ένα διερευνητικό χαρακτήρα, οπότε μπορούν να υπάρχουν πολλές βελτιώσεις και πεδία που μπορεί να εστιάσει κάποια μελλοντική έρευνα.
Σε επίπεδο βελτίωσης της αποδοτικότητας των αισθητήρων αέρα φθηνού κόστους που μπορούν να χρησιμοποιηθούν, μαζί με τον BME680 αλλά όχι μόνο, 
έχει μεγάλη σημασία να εφαρμόσουμε αυτό που προτείνετε σε έρευνες και ανασκοπήσεις \cite{CHOJER2020138385},\cite{SA2022102551}.
Συγκεκριμένα, υπάρχει η ανάγκη βαθμονόμησης του αισθητήρα σε εργαστηριακές συνθήκες αλλά και στον χώρο που θα χρησιμοποιηθούν με την χρήση κάποιου οργάνου αναφοράς.
Αυτό προτείνεται να γίνει και κατά την διάρκεια του χρόνου που βρίσκεται ένας αισθητήρας σε λειτουργία, σε μία περίοδο κάποιων μηνών, ώστε να αντιμετωπιστεί το sensor drift.

Για τους αισθητήρες αέρα φθηνού κόστους που καταμετράνε TVOCs, που είναι και ο BME680,έχει αξία να μπορεί να ελεγχθεί εργαστηριακά αν μπορεί να ξεχωρίσει κάποια πρωτογενή πηγή ρύπου που ξεφεύγει παραπάνω από το κανονικό και πως μπορούμε να ξεχωρίσουμε
από τις διαφορετικές πηγές TVOC.
Επίσης, παραπάνω αισθητήρες μπορούν να βοηθήσουν στην καλύτερη λειτουργία του τριπλευρικού εντοπισμού, καθώς αισθητήρες με χειρότερα αποτελέσματα θα επηρεάζουν λιγότερο το τελικό αποτέλεσμα.

Για τον BME680 ,μπορούν να χρησιμοποιηθούν μπορούν να χρησιμοποιηθούν και οι άλλες εξόδους τις BSEC, όπως το IAQ index και το CO2 equivalent, αλλά και οι περιβαλλοντικές τιμές που καταγράφει, όπως η υγρασία, η θερμοκρασία και η πίεση.
Θα μπορούσε να υπάρχει εργαστηριακή προσέγγιση για την καλύτερη γνώση των εσωτερικών μοντέλων τις BSEC ,αλλά δεν μπορούμε να πούμε με σιγουριά ότι αυτό θα ήταν αναγκαία.

Στο κομμάτι των πειραμάτων, η χρονική περίοδος για τα πειράματα επιλέχθηκε αυθαίρετα, αφότου δεν υπήρχε επαρκής γνώσης για την συμπεριφορά της εξατμισμένης αιθανόλης, λόγω έλλειψης γνώσης.
Η παραπάνω ειδικευμένη γνώση χημείας και φυσικής γνώσης αναμφίβολα θα βοηθήσει στην καλύτερη συγκέντρωση πληροφοριών και προσέγγιση μελλοντικών πειραμάτων.

Μπορεί να υπάρχει μεγάλη βελτίωση στο κομμάτι των μηχανικών μοντέλων, με χρήση πιο εξιδεικευμένων μοντέλων, διαφορετικές προσεγγίσεις εξαγωγής και επιλογής χαρακτηριστικών, όπως και χρήση βαθιάς μηχανικής μάθησης.
